\documentclass[12pt,a4paper]{article}

% XeLaTeX packages
\usepackage{fontspec}
\usepackage{polyglossia}
\usepackage{xcolor}
\usepackage{listings}
\usepackage{geometry}
\usepackage{titlesec}
\usepackage{mdframed}
\usepackage{tikz}
\usepackage{array}
\usepackage{booktabs}
\usepackage{enumitem}
\usepackage{fancyhdr}
\usepackage{hyperref}
\usepackage{tcolorbox}
\tcbuselibrary{listings,skins,breakable}

% Page geometry
\geometry{
  a4paper,
  top=2.5cm,
  bottom=2.5cm,
  left=2.5cm,
  right=2.5cm
}

% Fonts - using safe defaults for XeLaTeX
\setmainfont{DejaVu Serif}
\setmonofont{DejaVu Sans Mono}
\setsansfont{DejaVu Sans}

% Language
\setmainlanguage{english}

% Colors
\definecolor{arduinoBlue}{RGB}{0, 114, 178}
\definecolor{arduinoTeal}{RGB}{0, 158, 115}
\definecolor{codeBackground}{RGB}{245, 245, 245}
\definecolor{codeBorder}{RGB}{0, 114, 178}
\definecolor{commentColor}{RGB}{0, 128, 0}
\definecolor{keywordColor}{RGB}{0, 0, 200}
\definecolor{stringColor}{RGB}{180, 0, 0}
\definecolor{noteBackground}{RGB}{255, 255, 200}
\definecolor{warningBackground}{RGB}{255, 220, 200}
\definecolor{tipBackground}{RGB}{200, 240, 200}
\definecolor{headerColor}{RGB}{0, 80, 140}

% Listings style for Arduino/C++
\lstdefinestyle{arduino}{
  language=C++,
  backgroundcolor=\color{codeBackground},
  basicstyle=\ttfamily\small,
  keywordstyle=\color{keywordColor}\bfseries,
  commentstyle=\color{commentColor}\itshape,
  stringstyle=\color{stringColor},
  numberstyle=\tiny\color{gray},
  numbers=left,
  stepnumber=1,
  numbersep=8pt,
  frame=single,
  rulecolor=\color{codeBorder},
  breaklines=true,
  breakatwhitespace=false,
  showspaces=false,
  showstringspaces=false,
  tabsize=2,
  captionpos=b,
  morekeywords={
    pinMode, digitalWrite, digitalRead, analogWrite, analogRead,
    delay, delayMicroseconds, millis, micros,
    Serial, begin, print, println, available, read,
    LED_BUILTIN, INPUT, OUTPUT, INPUT_PULLUP,
    HIGH, LOW, true, false, void,
    setup, loop, int, bool, float, char, String,
    if, else, for, while, do, return,
    attachInterrupt, detachInterrupt, RISING, FALLING, CHANGE,
    tone, noTone, pulseIn, shiftOut, shiftIn,
    Wire, Servo, write, attach
  }
}

\lstset{style=arduino}

% tcolorbox styles
\tcbset{
  codebox/.style={
    colback=codeBackground,
    colframe=arduinoBlue,
    fonttitle=\bfseries\sffamily,
    title=#1,
    breakable,
    enhanced
  },
  notebox/.style={
    colback=noteBackground,
    colframe=yellow!60!black,
    fonttitle=\bfseries\sffamily,
    title=Yaad Rakho (Remember),
    breakable
  },
  tipbox/.style={
    colback=tipBackground,
    colframe=arduinoTeal,
    fonttitle=\bfseries\sffamily,
    title=Tip,
    breakable
  },
  warnbox/.style={
    colback=warningBackground,
    colframe=red!70!black,
    fonttitle=\bfseries\sffamily,
    title=Dhyan Do (Warning),
    breakable
  }
}

% Section formatting
\titleformat{\section}
  {\Large\bfseries\sffamily\color{headerColor}}
  {\thesection}{1em}{}[\titlerule]

\titleformat{\subsection}
  {\large\bfseries\sffamily\color{arduinoBlue}}
  {\thesubsection}{1em}{}

\titleformat{\subsubsection}
  {\normalsize\bfseries\sffamily\color{arduinoTeal}}
  {\thesubsubsection}{1em}{}

% Header/Footer
\pagestyle{fancy}
\fancyhf{}
\fancyhead[L]{\small\sffamily\color{arduinoBlue}Arduino C++ Notes}
\fancyhead[R]{\small\sffamily\color{arduinoBlue}Complete Guide}
\fancyfoot[C]{\small\thepage}
\renewcommand{\headrulewidth}{0.4pt}
\headrulecolor{arduinoBlue}

% Custom commands
\newcommand{\keyword}[1]{\textcolor{keywordColor}{\texttt{\textbf{#1}}}}
\newcommand{\func}[1]{\textcolor{arduinoBlue}{\texttt{#1()}}}
\newcommand{\pin}[1]{\textcolor{arduinoTeal}{\textbf{Pin #1}}}

\hypersetup{
  colorlinks=true,
  linkcolor=arduinoBlue,
  urlcolor=arduinoTeal
}

% ============================================================
\begin{document}
% ============================================================

% Title Page
\begin{titlepage}
  \centering
  \vspace*{2cm}
  
  {\Huge\bfseries\sffamily\color{arduinoBlue} Arduino C++ Notes}\\[0.5cm]
  {\Large\sffamily\color{arduinoTeal} Complete Programming Guide}\\[0.3cm]
  {\large\sffamily\color{gray} Beginner se Advanced tak}\\[2cm]
  
  \begin{tcolorbox}[colback=arduinoBlue!10, colframe=arduinoBlue, width=0.8\textwidth]
    \centering
    \large\sffamily
    Is notebook mein Arduino ke liye C++ ke\\
    sabhi important concepts aur examples hain.\\[0.3cm]
    \textbf{Topics:} Digital I/O $\bullet$ Analog $\bullet$ Serial $\bullet$\\
    Sensors $\bullet$ Motors $\bullet$ Interrupts $\bullet$ Timing
  \end{tcolorbox}
  
  \vfill
  {\large\sffamily\color{gray} Version 1.0}
\end{titlepage}

\tableofcontents
\newpage

% ============================================================
\section{Arduino Basics — Shuruat}
% ============================================================

\subsection{Arduino Program ki Structure}

Har Arduino program mein \textbf{do main functions} hote hain:

\begin{center}
\begin{tabular}{|l|l|l|}
\hline
\rowcolor{arduinoBlue!20}
\textbf{Function} & \textbf{Kab Chalta Hai} & \textbf{Kaam} \\
\hline
\texttt{setup()} & Sirf \textbf{ek baar} — start par & Sab kuch initialize karna \\
\hline
\texttt{loop()} & \textbf{Baar baar} — hamesha & Main kaam karna \\
\hline
\end{tabular}
\end{center}

\begin{tcolorbox}[codebox=Ek Basic Arduino Program]
\begin{lstlisting}
void setup() {
  // Yeh sirf ek baar chalega jab Arduino start hoga
  pinMode(LED_BUILTIN, OUTPUT); // Pin 13 ko output banao
}

void loop() {
  // Yeh baar baar chalta rahega hamesha
  digitalWrite(LED_BUILTIN, HIGH); // LED ON
  delay(1000);                      // 1 second ruko
  digitalWrite(LED_BUILTIN, LOW);  // LED OFF
  delay(1000);                      // 1 second ruko
}
\end{lstlisting}
\end{tcolorbox}

\subsubsection{Keywords ka Matlab}

\begin{itemize}[leftmargin=*]
  \item \keyword{void} — Function kuch return nahi karta (khali haath wapas)
  \item \keyword{setup} — Function ka naam (Arduino ka special naam)
  \item \texttt{()} — Function ka darwaza, andar parameters jaate hain
  \item \texttt{\{\}} — Function ka body, andar sab kaam likhte hain
  \item \texttt{;} — Har statement ke end mein lagta hai (zaroori!)
\end{itemize}

\begin{tcolorbox}[notebox]
\texttt{setup()} aur \texttt{loop()} Arduino mein \textbf{reserved names} hain. Inhe change mat karo — Arduino inhe automatically dhundta hai!
\end{tcolorbox}

\newpage

% ============================================================
\section{Digital Input / Output}
% ============================================================

\subsection{LED Blink — Simple On/Off}

\subsubsection{Concept}
\begin{itemize}
  \item \keyword{pinMode(pin, mode)} — Pin ko INPUT ya OUTPUT set karo
  \item \keyword{digitalWrite(pin, value)} — Pin ko HIGH (5V) ya LOW (0V) karo
  \item \keyword{delay(ms)} — Milliseconds tak ruko
\end{itemize}

\begin{tcolorbox}[codebox=Example 1: LED Blink]
\begin{lstlisting}
// Pin 13 par LED lagaao
int ledPin = 13; // Variable — pin number store karta hai

void setup() {
  pinMode(ledPin, OUTPUT); // ledPin ko output banao
}

void loop() {
  digitalWrite(ledPin, HIGH); // LED ON karo
  delay(500);                  // 0.5 second ruko
  digitalWrite(ledPin, LOW);  // LED OFF karo
  delay(500);                  // 0.5 second ruko
  // Yeh loop repeat hoga — LED blink karega
}
\end{lstlisting}
\end{tcolorbox}

\subsection{Multiple LEDs}

\begin{tcolorbox}[codebox=Example 2: Teen LED ka Chase Effect]
\begin{lstlisting}
// Teen alag pins par teen LED
int led1 = 11;
int led2 = 12;
int led3 = 13;

void setup() {
  pinMode(led1, OUTPUT);
  pinMode(led2, OUTPUT);
  pinMode(led3, OUTPUT);
}

void loop() {
  // Ek ek karke LED on karo — chase effect
  digitalWrite(led1, HIGH);
  delay(200);
  digitalWrite(led1, LOW);

  digitalWrite(led2, HIGH);
  delay(200);
  digitalWrite(led2, LOW);

  digitalWrite(led3, HIGH);
  delay(200);
  digitalWrite(led3, LOW);
}
\end{lstlisting}
\end{tcolorbox}

\subsection{Button se LED Control}

\subsubsection{digitalRead — Button Press Padhna}

\begin{tcolorbox}[codebox=Example 3: Button se LED On/Off]
\begin{lstlisting}
int buttonPin = 2;  // Button pin 2 par
int ledPin    = 13; // LED pin 13 par
int buttonState;    // Variable — button ki state store karta hai

void setup() {
  pinMode(buttonPin, INPUT);  // Button = INPUT
  pinMode(ledPin, OUTPUT);    // LED = OUTPUT
}

void loop() {
  buttonState = digitalRead(buttonPin); // Button padhо

  if (buttonState == HIGH) {      // Agar button dabaya
    digitalWrite(ledPin, HIGH);   // LED ON
  } else {                        // Warna
    digitalWrite(ledPin, LOW);    // LED OFF
  }
}
\end{lstlisting}
\end{tcolorbox}

\begin{tcolorbox}[tipbox]
\textbf{INPUT\_PULLUP} use karo agar external resistor nahi hai:\\
\texttt{pinMode(buttonPin, INPUT\_PULLUP);}\\
Is case mein logic ulta hoga — button dabane par \texttt{LOW} aayega.
\end{tcolorbox}

\begin{tcolorbox}[codebox=Example 4: Toggle Button (Press karo ON — Phir Press karo OFF)]
\begin{lstlisting}
int buttonPin = 2;
int ledPin    = 13;
bool ledState = false;    // LED abhi OFF hai
bool lastButton = false;  // Pehle button ka state

void setup() {
  pinMode(buttonPin, INPUT_PULLUP);
  pinMode(ledPin, OUTPUT);
}

void loop() {
  bool currentButton = (digitalRead(buttonPin) == LOW); // Dabaya?

  // Sirf tab toggle karo jab nayi press ho
  if (currentButton == true && lastButton == false) {
    ledState = !ledState;              // ON hoga toh OFF, OFF hoga toh ON
    digitalWrite(ledPin, ledState);   // LED update karo
  }

  lastButton = currentButton;
  delay(50); // Debounce — bounce se bachao
}
\end{lstlisting}
\end{tcolorbox}

\newpage

% ============================================================
\section{Analog Input / Output}
% ============================================================

\subsection{Analog Read — Sensor Value Padhna}

Arduino mein \textbf{Analog pins (A0-A5)} 0 se 1023 tak value dete hain.

\begin{center}
\begin{tabular}{|c|c|}
\hline
\rowcolor{arduinoBlue!20}
\textbf{Voltage} & \textbf{analogRead() Value} \\
\hline
0V & 0 \\
\hline
2.5V & 512 \\
\hline
5V & 1023 \\
\hline
\end{tabular}
\end{center}

\begin{tcolorbox}[codebox=Example 5: Potentiometer (Variable Resistor) se LED Control]
\begin{lstlisting}
int potPin = A0;   // Potentiometer A0 par lagaao
int ledPin = 9;    // LED PWM pin par (3,5,6,9,10,11)
int potValue;      // Potentiometer ki value (0-1023)
int brightness;    // LED ki brightness (0-255)

void setup() {
  pinMode(ledPin, OUTPUT);
  Serial.begin(9600); // Serial monitor shuru karo
}

void loop() {
  potValue   = analogRead(potPin);      // 0 se 1023 tak padhо
  brightness = potValue / 4;            // 0-1023 ko 0-255 mein badlo
  analogWrite(ledPin, brightness);      // LED ki brightness set karo

  Serial.print("Pot Value: ");
  Serial.print(potValue);
  Serial.print(" | Brightness: ");
  Serial.println(brightness);

  delay(100);
}
\end{lstlisting}
\end{tcolorbox}

\begin{tcolorbox}[notebox]
\textbf{analogWrite()} sirf PWM pins par kaam karta hai:\\
Arduino Uno mein: Pin \textbf{3, 5, 6, 9, 10, 11} (tilde \textasciitilde{} sign se pehchano)
\end{tcolorbox}

\subsection{map() Function — Value Range Badalna}

\begin{tcolorbox}[codebox=Example 6: map() se Temperature Display]
\begin{lstlisting}
int tempPin = A1;  // Temperature sensor A1 par

void setup() {
  Serial.begin(9600);
}

void loop() {
  int rawValue = analogRead(tempPin); // 0-1023

  // map(value, fromLow, fromHigh, toLow, toHigh)
  // 0-1023 ko 0-100 degree mein badlo
  int temperature = map(rawValue, 0, 1023, 0, 100);

  Serial.print("Temperature: ");
  Serial.print(temperature);
  Serial.println(" C");

  delay(500);
}
\end{lstlisting}
\end{tcolorbox}

\newpage

% ============================================================
\section{Serial Communication}
% ============================================================

\subsection{Serial Monitor se Baat Karna}

\begin{tcolorbox}[codebox=Example 7: Serial Monitor par Data Print Karna]
\begin{lstlisting}
int sensorPin = A0;

void setup() {
  Serial.begin(9600); // 9600 baud rate par serial shuru karo
  Serial.println("Arduino Ready!"); // Ek baar print karo
}

void loop() {
  int value = analogRead(sensorPin);

  Serial.print("Value = ");   // Bina newline ke print karo
  Serial.println(value);      // Value print karo + newline

  delay(1000);
}
\end{lstlisting}
\end{tcolorbox}

\subsection{Serial se Input Lena}

\begin{tcolorbox}[codebox=Example 8: Computer se LED Control karo]
\begin{lstlisting}
int ledPin = 13;
char command; // Ek character store karega

void setup() {
  pinMode(ledPin, OUTPUT);
  Serial.begin(9600);
  Serial.println("Commands: 'H' = ON, 'L' = OFF");
}

void loop() {
  if (Serial.available() > 0) { // Koi data aaya hai?
    command = Serial.read();    // Ek character padhо

    if (command == 'H') {
      digitalWrite(ledPin, HIGH);
      Serial.println("LED ON!");
    }
    else if (command == 'L') {
      digitalWrite(ledPin, LOW);
      Serial.println("LED OFF!");
    }
    else {
      Serial.println("Invalid command!");
    }
  }
}
\end{lstlisting}
\end{tcolorbox}

\newpage

% ============================================================
\section{Variables aur Data Types}
% ============================================================

\subsection{C++ mein Data Types}

\begin{center}
\begin{tabular}{|l|l|l|l|}
\hline
\rowcolor{arduinoBlue!20}
\textbf{Type} & \textbf{Size} & \textbf{Range} & \textbf{Example} \\
\hline
\texttt{int} & 2 bytes & -32768 to 32767 & \texttt{int x = 10;} \\
\hline
\texttt{long} & 4 bytes & -2B to 2B & \texttt{long t = 100000;} \\
\hline
\texttt{float} & 4 bytes & Decimal numbers & \texttt{float v = 3.14;} \\
\hline
\texttt{bool} & 1 bit & true / false & \texttt{bool on = true;} \\
\hline
\texttt{char} & 1 byte & One character & \texttt{char c = 'A';} \\
\hline
\texttt{String} & Variable & Text & \texttt{String s = "Hello";} \\
\hline
\texttt{byte} & 1 byte & 0 to 255 & \texttt{byte b = 255;} \\
\hline
\end{tabular}
\end{center}

\begin{tcolorbox}[codebox=Example 9: Variables ka Use]
\begin{lstlisting}
int    counter = 0;      // Poora number — counter ke liye
float  voltage = 0.0;    // Decimal number — voltage ke liye
bool   isOn    = false;  // Sirf true ya false
char   grade   = 'A';    // Ek akshar
String name    = "LED";  // Text
long   timer   = 0;      // Bada number — time ke liye

void setup() {
  Serial.begin(9600);
}

void loop() {
  counter++;  // counter = counter + 1 (shortcut)

  int raw = analogRead(A0);
  voltage = raw * (5.0 / 1023.0); // 0-1023 ko voltage mein badlo

  Serial.print("Count: "); Serial.println(counter);
  Serial.print("Voltage: "); Serial.println(voltage);

  delay(1000);
}
\end{lstlisting}
\end{tcolorbox}

\subsection{Arrays — Kai Values Ek Saath}

\begin{tcolorbox}[codebox=Example 10: Array se Multiple LEDs Control]
\begin{lstlisting}
// Array — kai values ek group mein
int ledPins[] = {8, 9, 10, 11, 12, 13}; // 6 LEDs ke pins
int numLeds   = 6;                        // Kitne LEDs hain

void setup() {
  // for loop se sabhi pins setup karo
  for (int i = 0; i < numLeds; i++) {
    pinMode(ledPins[i], OUTPUT); // ledPins[0], [1], [2]...
  }
}

void loop() {
  // Ek ek karke ON karo
  for (int i = 0; i < numLeds; i++) {
    digitalWrite(ledPins[i], HIGH);
    delay(100);
  }

  // Ek ek karke OFF karo
  for (int i = numLeds - 1; i >= 0; i--) {
    digitalWrite(ledPins[i], LOW);
    delay(100);
  }
}
\end{lstlisting}
\end{tcolorbox}

\newpage

% ============================================================
\section{Control Structures — Decision aur Loops}
% ============================================================

\subsection{if / else if / else}

\begin{tcolorbox}[codebox=Example 11: Light Sensor se Automatic LED]
\begin{lstlisting}
int ldrPin  = A0; // LDR (Light sensor) A0 par
int ledPin  = 13;
int ldrValue;

void setup() {
  pinMode(ledPin, OUTPUT);
  Serial.begin(9600);
}

void loop() {
  ldrValue = analogRead(ldrPin); // 0 = andhera, 1023 = ujala

  if (ldrValue < 300) {
    // Bahut andhera
    digitalWrite(ledPin, HIGH);
    Serial.println("Dark - LED ON");
  }
  else if (ldrValue < 700) {
    // Theek thak roshni
    Serial.println("Normal light");
  }
  else {
    // Bahut roshan
    digitalWrite(ledPin, LOW);
    Serial.println("Bright - LED OFF");
  }

  delay(500);
}
\end{lstlisting}
\end{tcolorbox}

\subsection{for Loop}

\begin{tcolorbox}[codebox=Example 12: for Loop se LED Fade In/Out]
\begin{lstlisting}
int ledPin = 9; // PWM pin zaroori hai

void setup() {
  pinMode(ledPin, OUTPUT);
}

void loop() {
  // Fade IN: 0 se 255 tak brightness badhao
  for (int brightness = 0; brightness <= 255; brightness++) {
    analogWrite(ledPin, brightness);
    delay(10); // Thoda ruko har step par
  }

  // Fade OUT: 255 se 0 tak brightness ghataо
  for (int brightness = 255; brightness >= 0; brightness--) {
    analogWrite(ledPin, brightness);
    delay(10);
  }
}
\end{lstlisting}
\end{tcolorbox}

\subsection{while Loop}

\begin{tcolorbox}[codebox=Example 13: while Loop — Button Dabne tak Wait Karo]
\begin{lstlisting}
int buttonPin = 2;
int ledPin    = 13;

void setup() {
  pinMode(buttonPin, INPUT_PULLUP);
  pinMode(ledPin, OUTPUT);
  Serial.begin(9600);
}

void loop() {
  Serial.println("Button dabao...");

  // Jab tak button nahi daba, yahan ruko
  while (digitalRead(buttonPin) == HIGH) {
    // Kuch mat karo — bas wait karo
    delay(10);
  }

  // Button dab gaya!
  Serial.println("Button daba gaya!");
  digitalWrite(ledPin, HIGH);
  delay(1000);
  digitalWrite(ledPin, LOW);
}
\end{lstlisting}
\end{tcolorbox}

\newpage

% ============================================================
\section{Functions Banana — Khud ka Function}
% ============================================================

\subsection{Custom Function Kaise Banate Hain}

\begin{tcolorbox}[codebox=Example 14: Custom Function — LED Blink Function]
\begin{lstlisting}
int ledPin = 13;

// Custom function banana:
// void = kuch return nahi
// blinkLED = function ka naam
// (int pin, int times, int waitTime) = parameters
void blinkLED(int pin, int times, int waitTime) {
  for (int i = 0; i < times; i++) {
    digitalWrite(pin, HIGH);
    delay(waitTime);
    digitalWrite(pin, LOW);
    delay(waitTime);
  }
}

// Value return karne wala function
float voltageRead(int pin) {
  int raw = analogRead(pin);
  float volts = raw * (5.0 / 1023.0);
  return volts; // Value wapas bhejo
}

void setup() {
  pinMode(ledPin, OUTPUT);
  Serial.begin(9600);
}

void loop() {
  blinkLED(ledPin, 3, 200); // 3 baar blink karo, 200ms gap

  float v = voltageRead(A0); // Voltage padhо
  Serial.print("Voltage: ");
  Serial.println(v);

  delay(2000);
}
\end{lstlisting}
\end{tcolorbox}

\begin{tcolorbox}[tipbox]
\textbf{Functions banane ke fayde:}
\begin{itemize}
  \item Code baar baar likhna nahi padta
  \item Code samajhna aasaan hota hai
  \item Ek jagah fix karo — sab jagah fix ho jaata hai
\end{itemize}
\end{tcolorbox}

\newpage

% ============================================================
\section{Timing — millis() ka Use}
% ============================================================

\subsection{delay() vs millis()}

\begin{tcolorbox}[warnbox]
\textbf{delay() ki problem:} Jab \texttt{delay()} chal raha hota hai, Arduino \textbf{koi kaam nahi kar sakta} — button press bhi nahi padh sakta! Isliye bade programs mein \texttt{millis()} use karte hain.
\end{tcolorbox}

\begin{tcolorbox}[codebox=Example 15: millis() se Non-Blocking Blink]
\begin{lstlisting}
int ledPin = 13;
bool ledState        = false; // LED ki current state
long previousMillis  = 0;     // Pehli baar ka time
long interval        = 1000;  // 1000ms = 1 second

void setup() {
  pinMode(ledPin, OUTPUT);
}

void loop() {
  long currentMillis = millis(); // Abhi kitna time hua (ms mein)

  // Kya interval guzar gaya?
  if (currentMillis - previousMillis >= interval) {
    previousMillis = currentMillis; // Time save karo

    ledState = !ledState;           // State toggle karo
    digitalWrite(ledPin, ledState); // LED update karo
  }

  // Yahan aur kaam bhi kar sakte hain bina ruke!
  // Button padhna, sensor check karna, etc.
}
\end{lstlisting}
\end{tcolorbox}

\begin{tcolorbox}[codebox=Example 16: millis() se Multiple Tasks Ek Saath]
\begin{lstlisting}
int led1 = 12, led2 = 13;
bool state1 = false, state2 = false;
long prev1 = 0, prev2 = 0;
long interval1 = 500, interval2 = 1300; // Alag alag speeds

void setup() {
  pinMode(led1, OUTPUT);
  pinMode(led2, OUTPUT);
}

void loop() {
  long now = millis();

  // LED 1 — 500ms par blink
  if (now - prev1 >= interval1) {
    prev1   = now;
    state1  = !state1;
    digitalWrite(led1, state1);
  }

  // LED 2 — 1300ms par blink (alag speed!)
  if (now - prev2 >= interval2) {
    prev2   = now;
    state2  = !state2;
    digitalWrite(led2, state2);
  }
  // Dono simultaneously blink karte hain — delay() se possible nahi tha!
}
\end{lstlisting}
\end{tcolorbox}

\newpage

% ============================================================
\section{Sensors ke Examples}
% ============================================================

\subsection{Ultrasonic Sensor — HC-SR04}

\begin{tcolorbox}[codebox=Example 17: Distance Measure karna]
\begin{lstlisting}
int trigPin = 9;  // Trigger pin
int echoPin = 10; // Echo pin

void setup() {
  pinMode(trigPin, OUTPUT);
  pinMode(echoPin, INPUT);
  Serial.begin(9600);
  Serial.println("Distance Sensor Ready");
}

float getDistance() {
  // Trigger pulse bhejo
  digitalWrite(trigPin, LOW);
  delayMicroseconds(2);
  digitalWrite(trigPin, HIGH);
  delayMicroseconds(10);
  digitalWrite(trigPin, LOW);

  // Echo ka time measure karo
  long duration = pulseIn(echoPin, HIGH);

  // Distance calculate karo (cm mein)
  // Sound ki speed = 343 m/s = 0.0343 cm/microsecond
  float distance = duration * 0.0343 / 2;
  return distance;
}

void loop() {
  float dist = getDistance();

  Serial.print("Distance: ");
  Serial.print(dist);
  Serial.println(" cm");

  if (dist < 10) {
    Serial.println("ALERT: Bahut paas!");
  }

  delay(500);
}
\end{lstlisting}
\end{tcolorbox}

\subsection{DHT11 — Temperature \& Humidity Sensor}

\begin{tcolorbox}[codebox=Example 18: Temp aur Humidity Read karna]
\begin{lstlisting}
// Library install karo: DHT sensor library by Adafruit
#include <DHT.h>

#define DHTPIN  2     // DHT11 pin 2 par
#define DHTTYPE DHT11 // DHT 11 ya DHT 22

DHT dht(DHTPIN, DHTTYPE); // DHT object banao

void setup() {
  Serial.begin(9600);
  dht.begin(); // Sensor shuru karo
}

void loop() {
  delay(2000); // DHT11 ko 2 seconds chahiye

  float humidity    = dht.readHumidity();    // Humidity %
  float temperature = dht.readTemperature(); // Celsius mein

  // Check karo reading sahi hai?
  if (isnan(humidity) || isnan(temperature)) {
    Serial.println("Sensor error!");
    return;
  }

  Serial.print("Humidity: ");
  Serial.print(humidity);
  Serial.print("% | Temperature: ");
  Serial.print(temperature);
  Serial.println(" C");
}
\end{lstlisting}
\end{tcolorbox}

\newpage

% ============================================================
\section{Motor Control}
% ============================================================

\subsection{Servo Motor}

\begin{tcolorbox}[codebox=Example 19: Servo Motor Control]
\begin{lstlisting}
#include <Servo.h> // Servo library

Servo myServo; // Servo object banao

void setup() {
  myServo.attach(9); // Servo ko pin 9 se attach karo
  Serial.begin(9600);
}

void loop() {
  // 0 se 180 degree tak sweep karo
  for (int angle = 0; angle <= 180; angle += 10) {
    myServo.write(angle); // Angle set karo
    Serial.print("Angle: ");
    Serial.println(angle);
    delay(100);
  }

  // 180 se 0 degree tak wapas aao
  for (int angle = 180; angle >= 0; angle -= 10) {
    myServo.write(angle);
    delay(100);
  }
}
\end{lstlisting}
\end{tcolorbox}

\subsection{DC Motor — L298N Motor Driver}

\begin{tcolorbox}[codebox=Example 20: DC Motor Forward/Backward]
\begin{lstlisting}
// L298N Motor Driver pins
int in1 = 8;  // Direction pin 1
int in2 = 7;  // Direction pin 2
int enA = 9;  // Enable pin (PWM — speed control)

void setup() {
  pinMode(in1, OUTPUT);
  pinMode(in2, OUTPUT);
  pinMode(enA, OUTPUT);
  Serial.begin(9600);
}

void motorForward(int speed) {
  digitalWrite(in1, HIGH); // Aage ki taraf
  digitalWrite(in2, LOW);
  analogWrite(enA, speed); // Speed 0-255
}

void motorBackward(int speed) {
  digitalWrite(in1, LOW);  // Peechhe ki taraf
  digitalWrite(in2, HIGH);
  analogWrite(enA, speed);
}

void motorStop() {
  digitalWrite(in1, LOW);
  digitalWrite(in2, LOW);
  analogWrite(enA, 0);
}

void loop() {
  Serial.println("Forward...");
  motorForward(200); // Full speed aage
  delay(2000);

  motorStop();
  delay(500);

  Serial.println("Backward...");
  motorBackward(150); // Slow speed peechhe
  delay(2000);

  motorStop();
  delay(1000);
}
\end{lstlisting}
\end{tcolorbox}

\newpage

% ============================================================
\section{Interrupts — Turant Response}
% ============================================================

\subsection{External Interrupt kya Hota Hai?}

Normal \texttt{loop()} mein code upar se neeche chalta hai — koi cheez bich mein rok nahi sakti. Interrupt ek \textbf{emergency signal} hai jo loop ko rokta hai aur turant ek function chalata hai.

\begin{tcolorbox}[codebox=Example 21: Button Interrupt — Instant Response]
\begin{lstlisting}
volatile bool ledState = false; // volatile zaroori hai interrupt mein

void ledToggle() { // Yeh function interrupt par chalega
  ledState = !ledState;
  digitalWrite(13, ledState);
}

void setup() {
  pinMode(13, OUTPUT);
  pinMode(2, INPUT_PULLUP); // Pin 2 interrupt ke liye

  // attachInterrupt(pin, function, when)
  // FALLING = jab HIGH se LOW ho (button dabane par)
  attachInterrupt(digitalPinToInterrupt(2), ledToggle, FALLING);

  Serial.begin(9600);
  Serial.println("Interrupt ready! Button dabao.");
}

void loop() {
  // Yahan koi bhi kaam karo
  // Button press par LED turant toggle hogi
  // loop() ko rokna nahi padega!
  Serial.println("Main loop chal raha hai...");
  delay(1000);
}
\end{lstlisting}
\end{tcolorbox}

\begin{tcolorbox}[notebox]
\textbf{Interrupt mein yaad rakho:}
\begin{itemize}
  \item Interrupt function mein \texttt{delay()} use mat karo
  \item Interrupt function chhota aur fast hona chahiye
  \item Shared variables ko \keyword{volatile} likhna zaroori hai
  \item Arduino Uno mein sirf Pin \textbf{2} aur \textbf{3} interrupt support karte hain
\end{itemize}
\end{tcolorbox}

\newpage

% ============================================================
\section{Quick Reference Card}
% ============================================================

\subsection{Sabse Common Functions}

\begin{center}
\begin{tabular}{|l|l|l|}
\hline
\rowcolor{arduinoBlue!20}
\textbf{Function} & \textbf{Parameters} & \textbf{Kaam} \\
\hline
\texttt{pinMode()} & pin, mode & Pin ka mode set karo \\
\hline
\texttt{digitalWrite()} & pin, value & Digital output bhejo \\
\hline
\texttt{digitalRead()} & pin & Digital input padhо \\
\hline
\texttt{analogWrite()} & pin, value(0-255) & PWM output bhejo \\
\hline
\texttt{analogRead()} & pin & Analog value padhо (0-1023) \\
\hline
\texttt{delay()} & milliseconds & Ruko \\
\hline
\texttt{millis()} & — & Ab tak ka time (ms) \\
\hline
\texttt{Serial.begin()} & baud rate & Serial shuru karo \\
\hline
\texttt{Serial.print()} & data & Print karo \\
\hline
\texttt{Serial.println()} & data & Print karo + newline \\
\hline
\texttt{map()} & val,fL,fH,tL,tH & Range badlo \\
\hline
\texttt{constrain()} & val, min, max & Range limit karo \\
\hline
\end{tabular}
\end{center}

\subsection{C++ Operators Cheat Sheet}

\begin{center}
\begin{tabular}{|l|l|l|}
\hline
\rowcolor{arduinoBlue!20}
\textbf{Operator} & \textbf{Naam} & \textbf{Example} \\
\hline
\texttt{=} & Assignment & \texttt{x = 5} \\
\hline
\texttt{==} & Equal check & \texttt{if (x == 5)} \\
\hline
\texttt{!=} & Not equal & \texttt{if (x != 0)} \\
\hline
\texttt{>} , \texttt{<} & Greater/Less & \texttt{if (x > 10)} \\
\hline
\texttt{++} & Increment & \texttt{x++} (x = x+1) \\
\hline
\texttt{--} & Decrement & \texttt{x--} (x = x-1) \\
\hline
\texttt{\&\&} & AND & \texttt{if (a \&\& b)} \\
\hline
\texttt{||} & OR & \texttt{if (a || b)} \\
\hline
\texttt{!} & NOT & \texttt{if (!a)} \\
\hline
\texttt{+=} & Add assign & \texttt{x += 5} (x = x+5) \\
\hline
\end{tabular}
\end{center}

\vspace{1cm}

\begin{tcolorbox}[colback=arduinoBlue!10, colframe=arduinoBlue, title=\textbf{All the Best!}]
\centering
\large\sffamily
Arduino seekhna ek exciting journey hai!\\
Chhote examples se shuru karo, experiment karo,\\
aur apne khud ke projects banao.\\[0.3cm]
\textbf{Har error ek naya seekhne ka mauka hai!}
\end{tcolorbox}

\end{document}
